\documentclass[12pt]{article}
\usepackage[utf8]{inputenc}
\usepackage[T1]{fontenc}
\usepackage[a4paper, left=2.4cm, right=2.4cm, top=2cm, bottom=2cm, nomarginpar]{geometry}
\usepackage{amsmath, algorithm, pgfplots, pgf, pdfpages, amsfonts, amssymb, stmaryrd, enumitem,graphicx, algpseudocode, setspace, listings, courier, color, caption, centernot, hyperref, lstautogobble, zi4, mathtools, lmodern, colortbl, array, multicol, fancyhdr, lastpage, svg}
\usepackage{booktabs}
\hypersetup{
    pdftitle={CE204_FP1_G6},
    pdftoolbar=true,        % show Acrobat’s toolbar?
    pdfmenubar=true,        % show Acrobat’s menu?
    pdfauthor={Muhammet Yağcıoğlu},
    pdfsubject={CE204_FP_RPs},
    pdfcreator={Muhammet Yağcıoğlu},   % creator of the document
    pdfproducer={Muhammet Yağcıoğlu}, % producer of the document
    colorlinks,
    linkcolor={black},
    citecolor={black},
    urlcolor={blue!80!black}
}
\usepackage[noblocks]{authblk}
\usepackage[version=4]{mhchem}
\usepackage[export]{adjustbox}
\usepackage{eso-pic} % Required for absolute positioning
\graphicspath{ {./images/} }
\usetikzlibrary{patterns}
\usepackage{tikz-3dplot}
\usepackage[version=4]{mhchem}
\usepackage[export]{adjustbox}
\usepgfplotslibrary{groupplots}
\usepackage{xcolor}
\usepackage[absolute,overlay]{textpos}
\usepackage{ragged2e}
\usepackage{xcolor}
\usepackage{amsthm}
\usepackage{tikz}




%----------compiler settings---------
\setlength{\parindent}{0pt}
\everymath{\displaystyle}
\onehalfspacing

%-----------colors----------
\pgfplotsset{compat=1.18}
\definecolor{bluekeywords}{rgb}{0.13, 0.13, 1}
\definecolor{greencomments}{rgb}{0, 0.5, 0}
\definecolor{redstrings}{rgb}{0.9, 0, 0}
\definecolor{graynumbers}{rgb}{0.5, 0.5, 0.5}
\definecolor{bronze}{rgb}{0.82, 0.41, 0.12}
\definecolor{crimson}{rgb}{0.86, 0.08, 0.24}
\definecolor{brown(web)}{rgb}{0.65, 0.16, 0.16}
\definecolor{darkspringgreen}{rgb}{0.16, 0.32, 0.75}
\definecolor{blue(ncs)}{rgb}{0.0, 0.53, 0.74}
\definecolor{amber}{rgb}{1.0, 0.49, 0.0}
\definecolor{cadmiumgreen}{rgb}{0.0, 0.42, 0.24}




%-------code frame-----------
\lstset{
    autogobble,
    columns=fullflexible,
    showspaces=false,
    showtabs=false,
    breaklines=true,
    showstringspaces=false,
    breakatwhitespace=true,
    escapeinside={(*@}{@*)},
    commentstyle=\color{greencomments},
    keywordstyle=\color{bluekeywords},
    stringstyle=\color{redstrings},
    numberstyle=\color{graynumbers},
    basicstyle=\ttfamily\footnotesize,
    frame=l,
    framesep=12pt,
    xleftmargin=12pt,
    tabsize=4,
    captionpos=b
}



\lstdefinelanguage{myMMA}{
keywords={SetDirectory, NotebookDirectory, Exp, IdentityMatrix, Eigenvalues, 
ListPlot, PlotRange, PlotStyle, Directive, PointSize, AspectRatio, Blue, Graphics, Line, 
Manipulate, Show, Sqrt, UniformDistribution, GammaDistribution, BetaDistribution, 
Nintegrate, For, DataRange, AxesLabel, PlotLabel, Transpose, Export, Plot, Append, Infinity},
keywordstyle=\color{black},
commentstyle=\color{gray}, 
identifierstyle=\color{blue},
sensitive=false,
comment=[l]{(*},
morecomment=[s]{/*}{*/},
morestring=[b]',
morestring=[b]"
}






%-----gauss declare---------
\pgfmathdeclarefunction{gauss}{2}{%
  \pgfmathparse{1/(#2*sqrt(2*pi)) * exp(-((x-#1)^2)/(2*#2^2))}%
}
\usetikzlibrary{arrows.meta, decorations.markings}
\pgfplotsset{compat=1.17}







%--------------page numbering------------
\fancyfoot[R]{page \thepage\ of \textcolor{bronze}{\pageref{LastPage}}}
\fancyhead[R]{Group 6}



%-------------qframe settings---------------

\usepackage[framemethod=TikZ]{mdframed}
\newmdenv[%
    skipabove=1em, % space above the frame
    skipbelow=1em, % space below the frame
    linewidth=1pt, % width of the frame lines
    linecolor=gray!80, % color of the frame lines
    backgroundcolor=gray!9, % background color inside the frame
    roundcorner=5pt, % radius of the rounded corners
    innerleftmargin=1em, % margin within the frame at the left
    innerrightmargin=1em, % margin within the frame at the right
    innertopmargin=0.5em, % margin within the frame at the top
    innerbottommargin=0.5em % margin within the frame at the bottom
]{qframe}

\newenvironment{q}
{
    \begin{qframe}
    \noindent\textit{\textbf{Problem Statement:}}
    \par\smallskip
}
{
    \end{qframe}
}





%-------------ANSWER TAG---------------
\newcommand{\AnswerTag}{\hfill \begin{tikzpicture} \fill[black] (0.35cm,0) -- (0,0.175cm) -- (0.35cm,0.35cm) -- cycle;\end{tikzpicture}}



%-------Theorem and another question frames----------
\newenvironment{theorem}{\begin{mdframed}\textbf{Theorem.}}{\end{mdframed}}

%usage: \begin{theorem}

\newenvironment{question}{\begin{mdframed}\textbf{Question.}}{\end{mdframed}}

%usage: %usage: \begin{question}
\fancyhead[L]{CE 224 Mechanics of Materials, Fall 2023, Homework-5 }
\title{\vspace{-1cm}CE 224 - HW5}
\author{Muhammet Yağcıoğlu - 290204042}
\usepackage{hyperref} 
\begin{document}
\maketitle\thispagestyle{fancy}
\pagestyle{fancy}
\tableofcontents
\newpage



\section*{Question 1}
\addcontentsline{toc}{section}{Question 1}
\begin{q}
A gondola on a ski lift is supported by two bent arms, as shown in the figure. Each arm is offset by the distance \(b=180 \mathrm{~mm}\) from the line of action of the weight force \(W\). The allowable stresses in the arms are \(100 \mathrm{MPa}\) in tension and \(50 \mathrm{MPa}\) in shear.

If the loaded gondola weighs \(12 \mathrm{kN}\), what is the minimum diameter \(d\) of the arms?
\end{q}

\begin{figure}[!ht]
    \centering
    \resizebox{0.7\textwidth}{!}{\includegraphics{Screenshot 2023-12-27 205128.png}}
    \caption{}
    \label{fig:enter-label}
\end{figure}


Let \(b\) be the offset distance of each arm, \(b = 180 \mathrm{~mm}\), and \(W\) the weight of the loaded gondola, \(W = 12 \mathrm{kN}\). The allowable tensile and shear stresses are denoted by \(\sigma_{\text{allow}} = 100 \mathrm{MPa}\) and \(\tau_{\text{allow}} = 50 \mathrm{MPa}\), respectively. Define \(d_{\text{min}}\) as the minimum diameter of the arms. Also note that \(\tau_{\max }=\frac{\sigma_t}{2}\). The cross-sectional area \(A\) and section modulus \(S\) are expressed as \(A = \frac{\pi d^2}{4}\) and \(S = \frac{\pi d^3}{32}\), respectively. For the determination of maximum tensile stress, employ: \[ \sigma_t = \frac{W}{A} + \frac{M}{S} = \frac{4W}{\pi d^2} + \frac{32Wb}{\pi d^3} \] or \[ \left(\frac{\pi \sigma_t}{4W}\right) d^3 - d - 8b = 0 \] Substituting numerical values: \[ \frac{\pi \sigma_t}{4W} = \frac{\pi \sigma_{\text{allow}}}{4W} = \frac{\pi \times 100 \mathrm{Mpa}}{4 \times 6 \mathrm{kN}} = 13,089.97 \frac{1}{\mathrm{~m}^2} \] and \(8b = 1.44 \mathrm{~m}\), the cubic equation becomes: \[ 13,090 d^3 - d - 1.44 = 0 \quad (d \text{ in meters}) \] Let \(c_1=13090\) and \(c_2=1.44\) for simplicity. Solve for \(d\) : \[c_1 d^3-d-c_2=0\] Divide both sides by \(c_1\) : \[ d^3-\frac{d}{c_1}-\frac{c_2}{c_1}=0 \] Change coordinates by substituting \(d=x+\frac{\lambda}{x}\), where \(\lambda\) is a constant value that will be determined later: \[ \left(x+\frac{\lambda}{x}\right)^3-\frac{x+\frac{\lambda}{x}}{c_1}-\frac{c_2}{c_1}=0 \] Multiply both sides by \(x^3\) and collect in terms of \(x\) : \[ x^6+x^4\left(3 \lambda-\frac{1}{c_1}\right)-\frac{x^3 c_2}{c_1}+x^2\left(3 \lambda^2-\frac{\lambda}{c_1}\right)+\lambda^3=0 \] Substitute \(\lambda=\frac{1}{3 c_1}\) and then \(y=x^3\), yielding a quadratic equation in the variable \(y\) : \[ y^2+\frac{1}{27 c_1^3}-\frac{y c_2}{c_1}=0 \] Find the positive solution to the quadratic equation: \[ y=\frac{9 c_2+\sqrt{3} c_1 \sqrt{\frac{27 c_1 c_2^2-4}{c_1^3}}}{18 c_1} \] Substitute back for \(y=x^3\) : \[ x^3=\frac{9 c_2+\sqrt{3} c_1 \sqrt{\frac{27 c_1 c_2^2-4}{c_1^3}}}{18 c_1} \] Taking cube roots gives \(\frac{\sqrt[3]{\frac{9 c_2+\sqrt{3} c_1 \sqrt{\frac{27 c_1 c_2^2-4}{c_1^3}}}{c_1}}}{\sqrt[3]{2} 3^{2 / \beta}}\) times the third roots of unity: \[x=\frac{\sqrt[3]{\frac{9 c_2+\sqrt{3} c_1 \sqrt{\frac{27 c_1 c_2^2-4}{c_1^3}}}{c_1}}}{\sqrt[3]{2} 3^{2 / 3}} \quad \text{or} \quad x=-\frac{\sqrt[3]{-\frac{1}{2}} \sqrt[3]{\frac{9 c_2+\sqrt{3} c_1 \sqrt{\frac{27 c_1 c_2^2-4}{c_1^3}}}{c_1}}}{3^{2 / 3}} \quad \text{or} \quad x=\frac{(-1)^{2 / 3} \sqrt[3]{\frac{9 c_2+\sqrt{3} c_1 \sqrt{\frac{27 c_1 c_2^2-4}{c_1^3}}}{c_1}}}{\sqrt[3]{2} 3^{2 / 3}}\] Substitute each value of \(x\) into \(d=x+\frac{1}{3 x c_1}\) and taking into real valued-d : \[d=\frac{\sqrt[3]{\frac{2}{3}}}{c_1 \sqrt[3]{\frac{9 c_2+\sqrt{3} c_1 \sqrt{\frac{27 c_1 c_2^2-4}{c_1^3}}}{c_1}}}+\frac{\sqrt[3]{\frac{9 c_2+\sqrt{3} c_1 \sqrt{\frac{27 c_1 c_2^2-4}{c_1^3}}}{c_1}}}{\sqrt[3]{2} 3^{2 / 3}}=\frac{\sqrt[3]{2} c_1\left(\frac{9 c_2}{c_1}+\sqrt{3} \sqrt{\frac{27 c_1 c_2^2-4}{c_1^3}}\right)^{2 / 3}+2 \sqrt[3]{3}}{6^{2 / 3} \sqrt[3]{\frac{9 c_2}{c_1}+\sqrt{3} \sqrt{\frac{27 c_1 c_2^2-4}{c_1^3}}} c_1}\] put back the constants \(c_1=13090\) and \(c_2=1.44\): \[d = \frac{2^{1/3} \cdot 13090 \left(\frac{9 \cdot 1.44}{13090} + \sqrt{3} \sqrt{\frac{27 \cdot 13090 \cdot (1.44)^2 - 4}{(13090)^3}}\right)^{2/3} + 2 \cdot 3^{1/3}}{6^{2/3} \sqrt[3]{\frac{9 \cdot 1.44}{13090} + \sqrt{3} \sqrt{\frac{27 \cdot 13090 \cdot (1.44)^2 - 4}{(13090)^3}}} \cdot 13090} \approx 0.0484467 \,\mathrm{~m}\] Hence \(d_{\text{min}} = 48.4 \mathrm{~mm}\).

\AnswerTag

\vfill
\begin{flushright}
\textbf{ans.} \(d_{\text{min}} = 48.4 \mathrm{~mm}\)
\end{flushright}

\newpage
\section*{Question 2}
\addcontentsline{toc}{section}{Question 2}
\begin{q}
2. The 20-kg drum is suspended from the hook mounted on the wooden frame. Determine the state of stress at point \(F\) on the cross section of the frame at section \(b-b\). Indicate the results on an element.
\end{q}

\begin{figure}[!ht]
    \centering
    \resizebox{0.4\textwidth}{!}{\includegraphics{Screenshot 2023-12-27 215441.png}}
    \caption{}
    \label{fig:enter-label}
\end{figure}

Let \(\Sigma\) represent the state of stress at point \(F\) on the cross-section of the wooden frame at section \(b-b\). The analysis begins with the determination of support reactions using equilibrium principles, referencing the free-body diagram of the entire frame:

\begin{figure}[!ht]
    \centering
    \includesvg[width=0.5\textwidth]{a2.svg}
    \caption{}
\end{figure}

The sum of the moments about point \(A\), denoted by \(\CircArrowLeft{+} \Sigma M_A = 0\), yields:

\[
F_{BD} \sin 30^{\circ}(3) - 20(9.81)(2) = 0 \implies F_{BD} = 261.6 \mathrm{~N}
\]

Next, considering the vertical equilibrium (\(+\uparrow \Sigma F_y = 0\)), we have:

\[
A_y - 261.6 \cos 30^{\circ} - 20(9.81) = 0 \implies A_y = 422.75 \mathrm{~N}
\]

The horizontal equilibrium (\(\stackrel{+}{\rightarrow} \Sigma F_x = 0\)) equation is:

\[
A_x - 261.6 \sin 30^{\circ} = 0 \implies A_x = 130.8 \mathrm{~N}
\]

Next, let us proceed to the analysis of internal loadings, delving into the equilibrium of the lower cut segment, as depicted in Fig. \(4\). At the right-angle corner, we invoke the equilibrium condition in the horizontal direction (\(+\uparrow \Sigma F_x = 0\)) to establish:

\[
130.8 - V = 0 \implies V = 130.8 \mathrm{~N}
\]

For the vertical equilibrium (\(+\uparrow \Sigma F_y = 0\)), we ascertain:

\[
422.75 - N = 0 \implies N = 422.75 \mathrm{~N}
\]

Additionally, by considering the sum of the moments about point \(K\) (\(\CircArrowLeft{+}\Sigma M_C = 0\)), the following equation emerges:

\[
130.8(1) - M = 0 \implies M = 130.8 \mathrm{~N \cdot m}
\]

\begin{figure}[!ht]
    \centering
    \includesvg[width=0.5\textwidth]{a3.svg}
    \caption{}
\end{figure}

Next, find the section properties, establishing the cross-sectional area \(A\) and moment of inertia \(I\). The cross-sectional area \(A\) is determined by multiplying the given dimensions, yielding:

\[
A = 0.075 \times 0.075 = 5.625 \times 10^{-3} \, \mathrm{m}^2
\]

Simultaneously, the moment of inertia \(I\) about the centroidal axis is computed using the standard formula for a rectangular section:

\[
I = \frac{1}{12} \times 0.075 \times (0.075)^3 = 2.6367 \times 10^{-6} \, \mathrm{m}^4
\]

Now, turning our attention to Fig. \(5\),


\begin{figure}[!ht]
    \centering
    \includesvg[width=0.4\textwidth]{a4.svg}
    \caption{}
\end{figure}


we evaluate the first moment of the area \(Q_F\) with respect to the neutral axis, employing the expression:

\[
Q_F = \bar{y}^{\prime} A^{\prime} = 0.025 \times 0.025 \times 0.075 = 46.875 \times 10^{-6} \, \mathrm{m}^3
\]

\textit{Normal stresses.} We commence with the foundational stress expression that captures the amalgamation of axial and bending stresses:

\[
\sigma = \frac{N}{A} \pm \frac{M y}{I}
\]

In our case, we focus on point \(F\) with a specific \(y\) value of \(0.0375 - 0.025 = 0.0125 \, \mathrm{m}\). The subsequent computation of the normal stress \(\sigma_F\) entails a series of precise steps:

\[ \begin{aligned}
\sigma_F & = \frac{-422.75}{5.625 \times 10^{-3}} - \frac{130.8 \times 0.0125}{2.6367 \times 10^{-6}} \\
& = -695.24 \, \mathrm{kPa} = 695 \, \mathrm{kPa} \, (\mathrm{Compression})
\end{aligned} \]
\AnswerTag

This outcome unveils the normal stress at point \(F\), with the negative sign indicating a compressive stress. The magnitude, \(695 \, \mathrm{kPa}\), means a quantitative measure of the stress magnitude, offering a mechanical behavior. 


Shear Stress: The shear stress is contributed by transverse shear stress only. Thus,
\[
\tau_A=\frac{V Q_A}{I t}=\frac{130.8\left[46.875\left(10^{-6}\right)\right]}{2.6367\left(10^{-6}\right)(0.075)}=31.0 \mathrm{kPa}
\]

\AnswerTag

The state of stress at point \(F\) is represented on the element shown following.


\begin{figure}[!ht]
    \centering
    \includesvg[width=0.4\textwidth]{a5.svg}
    \caption{}
\end{figure}
\begin{flushright}
\textbf{ans.}
\end{flushright}
\end{document}